\section{Theory: why a fiber field can decouple hardening from frequency}\label{sec:theory}

\subsection{Two energy densities over one stiffness}
With the mass-normalized adapted mode, the Rayleigh quotient and
Proposition~\ref{prop:relaxed} write both design quantities as integrals of local energy
densities:
\begin{equation}
\omega_1^2=\int_\Omega e_K\,\dd\Omega,\quad e_K=\bm\kappa_1^{\!\top}\bm D\,\bm\kappa_1;\qquad
\gamma=\int_\Omega e_G\,\dd\Omega,\quad e_G=\tfrac12\bm\varepsilon_r^{\!\top}\bm A\,\bm\varepsilon_r .
\label{eq:densities}
\end{equation}
$J=\gamma/(\omega_1^2w_c^2)$ is their ratio up to the mode's mid-plate value. Both
densities read the \emph{same} rotated material \eqref{eq:qbar}, but through different
strain fields: $\bm\kappa_1$ (second derivatives of $w_1$, largest where the mode bends)
and $\bm\varepsilon_r$ (squared first derivatives minus the in-plane relief). For the simply
supported fundamental both peak at the center, but with different principal directions
(Sec.~\ref{sec:mechanism}). The following proposition shows exactly how
each density depends on the local fiber angle.

\subsection{Local Fourier structure}
\begin{proposition}[Angular dependence of a laminate energy density]\label{prop:fourier}
Let $\bm s$ be an engineering strain vector whose tensor has principal values
$e_{\max},e_{\min}$, ordered so that $|e_{\max}|\ge|e_{\min}|$, and let $\hat\varphi$ be the
direction of $e_{\max}$. Then for the single-ply stiffness \eqref{eq:qbar},
\begin{equation}
\bm s^{\!\top}\bar{\bm Q}(\theta)\bm s=
\underbrace{U_1(e_{\max}^2+e_{\min}^2)+2U_4e_{\max}e_{\min}}_{\text{isotropic}}
+U_2\,\Delta\cos2(\theta-\hat\varphi)+U_3\,\Sigma\cos4(\theta-\hat\varphi),
\label{eq:fourier}
\end{equation}
with $\Delta=e_{\max}^2-e_{\min}^2\ge0$ and $\Sigma=(e_{\max}-e_{\min})^2\ge0$.
\end{proposition}
\begin{proof}
Rotation invariance gives $\bm s^{\!\top}\bar{\bm Q}(\theta)\bm s=\bm s'^{\!\top}\bar{\bm
Q}(\theta-\hat\varphi)\bm s'$, with $\bm s'=[e_{\max},e_{\min},0]^{\!\top}$ the strain in
its principal frame. Let $\psi=\theta-\hat\varphi$. Using \eqref{eq:qbar},
$\bar Q_{11}=U_1+U_2\cos2\psi+U_3\cos4\psi$,
$\bar Q_{22}=U_1-U_2\cos2\psi+U_3\cos4\psi$ and $\bar Q_{12}=U_4-U_3\cos4\psi$, we get
$\bar Q_{11}e_{\max}^2+2\bar Q_{12}e_{\max}e_{\min}+\bar Q_{22}e_{\min}^2
=U_1(e_{\max}^2+e_{\min}^2)+2U_4e_{\max}e_{\min}+U_2(e_{\max}^2-e_{\min}^2)\cos2\psi
+U_3(e_{\max}-e_{\min})^2\cos4\psi$.
\end{proof}

Applying \eqref{eq:fourier} to the two densities, pointwise:
\begin{align}
e_K(\bm x,\theta)&=\tfrac{h^3}{12}\big[I_K+U_2\Delta_K\cos2(\theta-\varphi_K)+U_3\Sigma_K\cos4(\theta-\varphi_K)\big],
\label{eq:eK}\\
e_G(\bm x,\theta)&=\tfrac{h}{2}\big[I_G+U_2\Delta_G\cos2(\theta-\varphi_G)+U_3\Sigma_G\cos4(\theta-\varphi_G)\big],
\label{eq:eG}
\end{align}
where $(\Delta_K,\Sigma_K,\varphi_K)$ come from the principal curvatures of $w_1$ and
$(\Delta_G,\Sigma_G,\varphi_G)$ from the principal relaxed stretches. For CFRP
$U_2\approx\nUtwoUthree\,U_3$, so the $2\theta$ harmonic dominates. To leading order, then:
\begin{itemize}[itemsep=1pt]
\item $e_K$ is largest (the point contributes most to $\omega_1^2$) when the fiber lies
\emph{along} the principal curvature, $\theta=\varphi_K$;
\item $e_G$ is smallest (the point contributes least to $\gamma$) when the fiber lies
\emph{across} the principal stretch, $\theta=\varphi_G+\pi/2$.
\end{itemize}
Which way the frequency pulls depends on the sign of the multiplier $\Lambda$ of
Proposition~\ref{prop:pointwise} below, which weights $e_K$ against $e_G$. If the target
frequency lies above the one the stretch alone would select, $\Lambda>0$ and the fiber is
drawn toward the principal curvature, $\theta=\varphi_K$. If it lies below, $\Lambda<0$
and the fiber is pushed \emph{across} the principal curvature, $\theta=\varphi_K+\pi/2$,
the bending-compliant direction. We call $\varphi_G+\pi/2$ and the frequency-favored
direction the \emph{preferred directions}, and the angle between them,
$\chi(\bm x)\in[0,\pi/2]$, the \emph{conflict}. Where $\chi=0$, one fiber direction
serves both energies. Where $\chi$ is large, any single direction compromises.

\begin{corollary}[The unidirectional floor]\label{cor:udfloor}
Among all uniform symmetric laminates of the same thickness and mass, whatever the
stacking sequence, the lowest linear frequency is attained by a single-angle laminate.
\end{corollary}
\begin{proof}
For a symmetric laminate $\bm D=\tfrac{h^3}{12}\big(\bm\Gamma_0+\sum_k\xi^D_k\bm\Gamma_k\big)$ is affine
in its four lamination parameters $\bm\xi^D$, and the set of $\bm\xi^D$ reachable by
stacking lies in the convex hull of the single-angle curve
$\bm\xi^D(\theta)=(\cos2\theta,\sin2\theta,\cos4\theta,\sin4\theta)$ \cite{tsai1968}.
$\omega_1^2$ is the minimum over admissible $w$ of a Rayleigh quotient that is affine in
$\bm\xi^D$, hence a concave function of $\bm\xi^D$, and a concave function on a convex
hull attains its minimum at an extreme point, i.e.\ at a single-angle laminate.
\end{proof}
\noindent For the plate studied here the single-angle frequency rises monotonically from
$\theta=0^\circ$ (Sec.~\ref{sec:levers}), so UD is the lowest-frequency uniform laminate of
\emph{any} stacking sequence, not only of the enumerated ones, and at $\omega_{1,\UD}$ there
is no other uniform competitor.

\begin{corollary}[P1: isotropy annihilates the lever]\label{cor:iso}
If $U_2=U_3=0$, then $\bm A$ and $\bm D$ are independent of every ply angle at every point.
Hence $\omega_1$, $\gamma$ and $J$ are independent of the orientation field, uniform or
steered or stacked, and every orientation sensitivity vanishes identically. This is exact,
including mode adaptation and relaxation, because the operators themselves do not depend
on $\theta$.
\end{corollary}
\noindent Numerically, $\partial J/\partial\theta\equiv0$ to $10^{-10}$ relative when
$\bm Q$ is set isotropic (\texttt{tests/test\_relaxed\_plate.py}).

\subsection{Two separation channels}\label{sec:channels}
Inserting \eqref{eq:lampar} into \eqref{eq:densities} gives
\begin{equation}
e_K=\tfrac{h^3}{12}\Big[\bm\kappa_1^{\!\top}\bm\Gamma_0\bm\kappa_1+\sum_i\xi^D_i(\bm x)\,
\bm\kappa_1^{\!\top}\bm\Gamma_i\bm\kappa_1\Big],\qquad
e_G=\tfrac{h}{2}\Big[\bm\varepsilon_r^{\!\top}\bm\Gamma_0\bm\varepsilon_r+\sum_i\xi^A_i(\bm x)\,
\bm\varepsilon_r^{\!\top}\bm\Gamma_i\bm\varepsilon_r\Big].
\label{eq:channels}
\end{equation}
\emph{The frequency reads only the bending lamination parameters, and the hardening only
the membrane ones.} Each design class corresponds to a different freedom in the field
$(\bm\xi^A(\bm x),\bm\xi^D(\bm x))$:

\begin{center}\small
\begin{tabular}{@{}llll@{}}
\toprule
class & $\bm\xi^A$ & $\bm\xi^D$ & separation available\\ \midrule
uniform single angle & $\bm v(\theta_0)$ & $=\bm\xi^A$ & none (one scalar)\\
uniform stack & const. & const.\ $\ne\bm\xi^A$ & through the thickness\\
steered single ply & $\bm v(\theta(\bm x))$ & $=\bm\xi^A$ & in the plane\\
steered stack & $\mathcal R_{\theta(\bm x)}\bm\xi^A_0$ & $\mathcal R_{\theta(\bm x)}\bm\xi^D_0$ & both\\
\bottomrule
\end{tabular}
\end{center}

\noindent Here $\mathcal R_\theta$ rotates the parameter pairs,
$(\xi_1+i\xi_2)\mapsto e^{2i\theta}(\xi_1+i\xi_2)$ and
$(\xi_3+i\xi_4)\mapsto e^{4i\theta}(\xi_3+i\xi_4)$. Section~\ref{sec:levers} measures the
achievable $(\omega_1,J)$ set of each class.

\subsection{Pointwise optimality: what an optimal field must do}\label{sec:pointwise}
Consider the problem
\begin{equation}
\min_{\theta(\cdot)}\;J\quad\text{s.t.}\quad\omega_1=\omega_\star
\label{eq:heldproblem}
\end{equation}
over single-ply fields. Freeze $w_1$ and $\bm\varepsilon_r$ at a reference design. This
\emph{frozen-field} model is exact to first order in the design change: $\bm
s^\star$ is stationary (Proposition~\ref{prop:relaxed}) and $\omega_1^2$ is a Rayleigh
quotient. In this model $w_c$ is fixed, so \eqref{eq:heldproblem} reads
$\min\int e_G(\bm x,\theta(\bm x))$ subject to $\int e_K(\bm x,\theta(\bm x))=\omega_\star^2$.

\begin{proposition}[Pointwise optimality and strict benefit of steering]\label{prop:pointwise}
In the frozen-field model the optimal field is
\begin{equation}
\theta^\star(\bm x)=\argmin_\theta\big[e_G(\bm x,\theta)-\Lambda\,e_K(\bm x,\theta)\big],
\label{eq:thetastar}
\end{equation}
with the scalar multiplier $\Lambda$ fixed by the frequency constraint. A uniform angle
$\theta_u$ can be optimal among all fields only if
\begin{equation}
U_2\Delta_G\sin2(\theta_u-\varphi_G)+2U_3\Sigma_G\sin4(\theta_u-\varphi_G)
=\Lambda'\big[U_2\Delta_K\sin2(\theta_u-\varphi_K)+2U_3\Sigma_K\sin4(\theta_u-\varphi_K)\big]
\label{eq:uniformkkt}
\end{equation}
holds at almost every $\bm x$, with one constant $\Lambda'=\tfrac{h^2}{6}\Lambda$.
Otherwise a steered field strictly lowers $J$ at the same frequency.
\end{proposition}
\begin{proof}
The Lagrangian $\int e_G-\Lambda(\int e_K-\omega_\star^2)$ has no derivative coupling
between points, so it is minimized pointwise, giving \eqref{eq:thetastar}. Every
uniform angle is a feasible field, so the field optimum is never worse. If the uniform
optimum were also the field optimum, it would satisfy the pointwise stationarity
condition $\partial_\theta e_G=\Lambda\partial_\theta e_K$ at every point. By
\eqref{eq:eK}--\eqref{eq:eG} this is \eqref{eq:uniformkkt}.
\end{proof}

Condition \eqref{eq:uniformkkt} asks the anisotropic parts of two different strain
fields to be proportional with one constant at a single angle. That fails whenever the
conflict $\chi(\bm x)$ or the ratio $\Delta_G/\Delta_K$ varies over the plate. Since the
principal curvature and the principal relaxed stretch of a bending mode turn differently
over the plate, and their energy ratio varies, it fails generically. Proposition~\ref{prop:pointwise} is the
precise form of the hypothesis. It also yields a test (P4): $\theta^\star$ computed from
the \emph{unidirectional} fields alone, with no optimizer, should already capture much
of the gain, and the optimized field should resemble it.

\subsection{One dimension in closed form}\label{sec:beamtheory}
\begin{proposition}[Beam: harmonic against curvature-weighted mean]\label{prop:beam}
Consider a symmetric single-ply beam with immovable ends, $A=bhE(x)$ and
$D=\tfrac{bh^3}{12}E(x)$, and a mass-normalized fundamental $\varphi$ with mid-span value
$\varphi_c$. Then
\begin{equation}
\gamma=\frac{bh\,(\int_0^L\varphi'^2)^2}{2\int_0^LE^{-1}},\qquad
\omega_1^2=\frac{bh^3}{12}\int_0^LE\,\varphi''^2,\qquad
J=\frac{6}{h^2\varphi_c^2}\,
\frac{(\int\varphi'^2)^2}{\int E^{-1}\;\int E\varphi''^2}.
\label{eq:beamJ}
\end{equation}
\end{proposition}
\begin{proof}
Axial equilibrium makes $N$ uniform. Integrating $\varepsilon=N/A=u'+\tfrac12q^2\varphi'^2$
over the span with $u(0)=u(L)=0$ gives $N\int A^{-1}=\tfrac12q^2\int\varphi'^2$. The
membrane energy is then $\tfrac12N^2\int A^{-1}=q^4(\int\varphi'^2)^2/(8\int A^{-1})$,
which equals $\tfrac\gamma4q^4$. The frequency is the Rayleigh quotient, and $J$ follows
from \eqref{eq:J}.
\end{proof}

\begin{corollary}[P2: uniform invariance]\label{cor:beaminv}
For uniform $E\equiv E_0$, $\int E^{-1}\int E\varphi''^2=L\int\varphi''^2$, and $\varphi$
itself is independent of $E_0$ (the operators scale together). Hence $J$ is
\emph{exactly} independent of any uniform single-ply angle, and every change in $J$ is
caused by spatial variation of $E(x)$.
\end{corollary}

Equation \eqref{eq:beamJ} makes the mechanism transparent. The hardening sees the
\emph{harmonic} mean of $E$ over the whole span, which is dominated by the softest
segment. The frequency sees the mean of $E$ \emph{weighted by the curvature} $\varphi''^2$,
which vanishes at the inflection points ($x/L\approx0.224$ and $0.776$ for the
clamped fundamental). Compliant fibers placed there raise $\int E^{-1}$, and so lower the
hardening, at almost no cost in frequency.

\begin{corollary}[Bang-bang bound]\label{cor:bang}
Suppose $E=E_{\max}$ except on zones of total length $\ell$ where $E=E_{\min}$ and
$\varphi''\approx0$. Then, to frozen-mode accuracy,
\begin{equation}
J/J_{\UD}\approx\big[1+(\ell/L)(E_{\max}/E_{\min}-1)\big]^{-1}.
\label{eq:bang}
\end{equation}
\end{corollary}
\noindent For CFRP, $E_{\max}/E_{\min}=\bar Q_{11}(0)/\bar Q_{11}(90^\circ)=14.6$, so
$\ell/L=0.05$ already gives $J/J_{\UD}\approx0.59$. Equation \eqref{eq:bang} contains P5
in one dimension. The lever grows with the anisotropy ratio and with how narrow the
compliant zones can be made while staying where $\varphi''\approx0$. The narrowness is set
by the tow curvature: turning a tow by $90^\circ$ at rate $\kappa_{\mathrm{tow}}$ needs a
run of $(\pi/2)/\kappa_{\mathrm{tow}}$, which is $\SI{0.10}{m}$ at
$900^\circ/\mathrm{m}$, a third of the span.
